\begin{helper}
For every $\eta>0$ there are real parameters $\alpha,L,\gamma_0$ and a
natural number $\Delta_0$ such that
\[
0<\alpha\le 1,\qquad 0<1-\alpha,\qquad 0\le L,\qquad 0\le\gamma_0,
\qquad 3\alpha\le \eta/3,
\]
\[
e^{-\gamma_0}\left(1+\frac1\alpha\right)\le\alpha,
\]
\[
\frac53+2\alpha\le
2(1-\alpha)^2 e^{-\alpha L}
\int_0^L\bigl(e^{-x}-e^{-L}\bigr)\,dx,
\]
and, for every natural $\Delta\ge\Delta_0$,
\[
\gamma_0<\Delta,\qquad L<\Delta,\qquad
\frac{\gamma_0^2}{\Delta-\gamma_0}\le-\log(1-\alpha),\qquad
\frac{L^2}{\Delta-L}\le-\log(1-\alpha).
\]
\end{helper}
\begin{proof}
Choose
\[
\alpha=\min\{10^{-3},\eta/9\}.
\]
Since $\eta>0$, this gives $0<\alpha\le10^{-3}\le1$ and
$3\alpha\le\eta/3$; in particular $0<1-\alpha$.  Put $L=6$.

The exponential tail tends to zero as $t\to\infty$, so we may choose a real
$g$ such that
\[
e^{-t}\left(1+\frac1\alpha\right)\le \alpha
\quad\text{for all }t\ge g .
\]
Set $\gamma_0=\max\{g,0\}$.  Then $0\le\gamma_0$, and the same tail
inequality holds at $\gamma_0$.

It remains to check the fixed numerical margin.  From
$2<e$, raising to the sixth power gives $64<e^6$, hence
$e^{-6}\le 1/64$.  Also
\[
\int_0^6 (e^{-x}-e^{-6})\,dx
=1-7e^{-6}\ge 57/64.
\]
Because $\alpha\le10^{-3}$, we have
\[
(1-\alpha)^2\ge (999/1000)^2
\quad\text{and}\quad
e^{-6\alpha}\ge 1-6\alpha\ge 497/500,
\]
where the exponential lower bound is the elementary inequality
$1+y\le e^y$ applied to $y=-6\alpha$.  Therefore
\[
2(1-\alpha)^2e^{-6\alpha}
\int_0^6(e^{-x}-e^{-6})\,dx
\ge
2\left(\frac{998001}{1000000}\right)
\left(\frac{497}{500}\right)
\left(\frac{57}{64}\right).
\]
The right hand side is greater than $5/3+2/1000$, and
$2\alpha\le2/1000$, so the desired high-overlap numerical inequality holds
for $L=6$.

Finally, since $0<\alpha<1$, we have $-\log(1-\alpha)>0$.  As
$\Delta\to\infty$ through the natural numbers, both $\Delta$ and
$\Delta-\gamma_0$ tend to infinity, and so do $\Delta-L$; hence
\[
\frac{\gamma_0^2}{\Delta-\gamma_0}\to0,
\qquad
\frac{L^2}{\Delta-L}\to0.
\]
Choosing $\Delta_0$ large enough gives simultaneously
$\gamma_0<\Delta$, $L<\Delta$, and the two displayed logarithmic error
bounds for every natural $\Delta\ge\Delta_0$.
\end{proof}
